\documentclass[11pt]{article}
\newcommand{\videoid}{HlZCNLpj8qw}
% Preámbulo compartido de las notas del curso de topología algebraica.
% Cada nota empieza con \documentclass[11pt]{article}, define \videoid{...}
% (el identificador del video de YouTube) y luego hace % Preámbulo compartido de las notas del curso de topología algebraica.
% Cada nota empieza con \documentclass[11pt]{article}, define \videoid{...}
% (el identificador del video de YouTube) y luego hace % Preámbulo compartido de las notas del curso de topología algebraica.
% Cada nota empieza con \documentclass[11pt]{article}, define \videoid{...}
% (el identificador del video de YouTube) y luego hace \input{preambulo}.
% Compilar con LuaLaTeX:  lualatex nombre-de-la-nota.tex
\usepackage[spanish,provide=*]{babel}
\usepackage[T1]{fontenc}
\usepackage{lmodern}
\usepackage[a4paper,margin=2.6cm]{geometry}
\usepackage{amsmath,amssymb,amsthm}
\usepackage{xcolor}
\usepackage{tikz}
\usetikzlibrary{arrows.meta,babel}
\usepackage[colorlinks=true,urlcolor=blue!50!black,linkcolor=blue!50!black]{hyperref}

\hyphenpenalty=10000 \exhyphenpenalty=10000 \emergencystretch=3em

% Colores de las figuras (los mismos que usa la página)
\definecolor{extremo}{RGB}{47,107,79}
\definecolor{azul}{RGB}{44,74,154}

% Definiciones numeradas por unidad
\theoremstyle{definition}
\newtheorem{defn}{Definición}
% Numeración por unidad: \unidad{u}{n} (después de \input{preambulo}) hace que las
% definiciones de la nota se llamen u.n, u.(n+1), ...; n es el primer número
% libre de esa unidad (ver la tabla de CONVENCIONES.md).
\newcommand{\unidad}[2]{\renewcommand{\thedefn}{#1.\arabic{defn}}\setcounter{defn}{\numexpr#2-1\relax}}

% Observaciones y erratas del video. Usar SOLO si hay algo que señalar:
%   \begin{erratas}
%     \errata{203}{3:23}{Se dice ..., pero debe decir ...}
%   \end{erratas}
% El primer argumento son los segundos (enlaza al minuto exacto del video),
% el segundo es la marca visible, el tercero es el texto.
\newenvironment{erratas}{\section*{Observaciones y erratas del video}\begin{itemize}}{\end{itemize}}
\newcommand{\errata}[3]{%
  \edef\errurl{https://www.youtube.com/watch?v=\videoid\string&t=#1}%
  \item[\textbf{\expandafter\href\expandafter{\errurl}{#2}}] #3}
.
% Compilar con LuaLaTeX:  lualatex nombre-de-la-nota.tex
\usepackage[spanish,provide=*]{babel}
\usepackage[T1]{fontenc}
\usepackage{lmodern}
\usepackage[a4paper,margin=2.6cm]{geometry}
\usepackage{amsmath,amssymb,amsthm}
\usepackage{xcolor}
\usepackage{tikz}
\usetikzlibrary{arrows.meta,babel}
\usepackage[colorlinks=true,urlcolor=blue!50!black,linkcolor=blue!50!black]{hyperref}

\hyphenpenalty=10000 \exhyphenpenalty=10000 \emergencystretch=3em

% Colores de las figuras (los mismos que usa la página)
\definecolor{extremo}{RGB}{47,107,79}
\definecolor{azul}{RGB}{44,74,154}

% Definiciones numeradas por unidad
\theoremstyle{definition}
\newtheorem{defn}{Definición}
% Numeración por unidad: \unidad{u}{n} (después de % Preámbulo compartido de las notas del curso de topología algebraica.
% Cada nota empieza con \documentclass[11pt]{article}, define \videoid{...}
% (el identificador del video de YouTube) y luego hace \input{preambulo}.
% Compilar con LuaLaTeX:  lualatex nombre-de-la-nota.tex
\usepackage[spanish,provide=*]{babel}
\usepackage[T1]{fontenc}
\usepackage{lmodern}
\usepackage[a4paper,margin=2.6cm]{geometry}
\usepackage{amsmath,amssymb,amsthm}
\usepackage{xcolor}
\usepackage{tikz}
\usetikzlibrary{arrows.meta,babel}
\usepackage[colorlinks=true,urlcolor=blue!50!black,linkcolor=blue!50!black]{hyperref}

\hyphenpenalty=10000 \exhyphenpenalty=10000 \emergencystretch=3em

% Colores de las figuras (los mismos que usa la página)
\definecolor{extremo}{RGB}{47,107,79}
\definecolor{azul}{RGB}{44,74,154}

% Definiciones numeradas por unidad
\theoremstyle{definition}
\newtheorem{defn}{Definición}
% Numeración por unidad: \unidad{u}{n} (después de \input{preambulo}) hace que las
% definiciones de la nota se llamen u.n, u.(n+1), ...; n es el primer número
% libre de esa unidad (ver la tabla de CONVENCIONES.md).
\newcommand{\unidad}[2]{\renewcommand{\thedefn}{#1.\arabic{defn}}\setcounter{defn}{\numexpr#2-1\relax}}

% Observaciones y erratas del video. Usar SOLO si hay algo que señalar:
%   \begin{erratas}
%     \errata{203}{3:23}{Se dice ..., pero debe decir ...}
%   \end{erratas}
% El primer argumento son los segundos (enlaza al minuto exacto del video),
% el segundo es la marca visible, el tercero es el texto.
\newenvironment{erratas}{\section*{Observaciones y erratas del video}\begin{itemize}}{\end{itemize}}
\newcommand{\errata}[3]{%
  \edef\errurl{https://www.youtube.com/watch?v=\videoid\string&t=#1}%
  \item[\textbf{\expandafter\href\expandafter{\errurl}{#2}}] #3}
) hace que las
% definiciones de la nota se llamen u.n, u.(n+1), ...; n es el primer número
% libre de esa unidad (ver la tabla de CONVENCIONES.md).
\newcommand{\unidad}[2]{\renewcommand{\thedefn}{#1.\arabic{defn}}\setcounter{defn}{\numexpr#2-1\relax}}

% Observaciones y erratas del video. Usar SOLO si hay algo que señalar:
%   \begin{erratas}
%     \errata{203}{3:23}{Se dice ..., pero debe decir ...}
%   \end{erratas}
% El primer argumento son los segundos (enlaza al minuto exacto del video),
% el segundo es la marca visible, el tercero es el texto.
\newenvironment{erratas}{\section*{Observaciones y erratas del video}\begin{itemize}}{\end{itemize}}
\newcommand{\errata}[3]{%
  \edef\errurl{https://www.youtube.com/watch?v=\videoid\string&t=#1}%
  \item[\textbf{\expandafter\href\expandafter{\errurl}{#2}}] #3}
.
% Compilar con LuaLaTeX:  lualatex nombre-de-la-nota.tex
\usepackage[spanish,provide=*]{babel}
\usepackage[T1]{fontenc}
\usepackage{lmodern}
\usepackage[a4paper,margin=2.6cm]{geometry}
\usepackage{amsmath,amssymb,amsthm}
\usepackage{xcolor}
\usepackage{tikz}
\usetikzlibrary{arrows.meta,babel}
\usepackage[colorlinks=true,urlcolor=blue!50!black,linkcolor=blue!50!black]{hyperref}

\hyphenpenalty=10000 \exhyphenpenalty=10000 \emergencystretch=3em

% Colores de las figuras (los mismos que usa la página)
\definecolor{extremo}{RGB}{47,107,79}
\definecolor{azul}{RGB}{44,74,154}

% Definiciones numeradas por unidad
\theoremstyle{definition}
\newtheorem{defn}{Definición}
% Numeración por unidad: \unidad{u}{n} (después de % Preámbulo compartido de las notas del curso de topología algebraica.
% Cada nota empieza con \documentclass[11pt]{article}, define \videoid{...}
% (el identificador del video de YouTube) y luego hace % Preámbulo compartido de las notas del curso de topología algebraica.
% Cada nota empieza con \documentclass[11pt]{article}, define \videoid{...}
% (el identificador del video de YouTube) y luego hace \input{preambulo}.
% Compilar con LuaLaTeX:  lualatex nombre-de-la-nota.tex
\usepackage[spanish,provide=*]{babel}
\usepackage[T1]{fontenc}
\usepackage{lmodern}
\usepackage[a4paper,margin=2.6cm]{geometry}
\usepackage{amsmath,amssymb,amsthm}
\usepackage{xcolor}
\usepackage{tikz}
\usetikzlibrary{arrows.meta,babel}
\usepackage[colorlinks=true,urlcolor=blue!50!black,linkcolor=blue!50!black]{hyperref}

\hyphenpenalty=10000 \exhyphenpenalty=10000 \emergencystretch=3em

% Colores de las figuras (los mismos que usa la página)
\definecolor{extremo}{RGB}{47,107,79}
\definecolor{azul}{RGB}{44,74,154}

% Definiciones numeradas por unidad
\theoremstyle{definition}
\newtheorem{defn}{Definición}
% Numeración por unidad: \unidad{u}{n} (después de \input{preambulo}) hace que las
% definiciones de la nota se llamen u.n, u.(n+1), ...; n es el primer número
% libre de esa unidad (ver la tabla de CONVENCIONES.md).
\newcommand{\unidad}[2]{\renewcommand{\thedefn}{#1.\arabic{defn}}\setcounter{defn}{\numexpr#2-1\relax}}

% Observaciones y erratas del video. Usar SOLO si hay algo que señalar:
%   \begin{erratas}
%     \errata{203}{3:23}{Se dice ..., pero debe decir ...}
%   \end{erratas}
% El primer argumento son los segundos (enlaza al minuto exacto del video),
% el segundo es la marca visible, el tercero es el texto.
\newenvironment{erratas}{\section*{Observaciones y erratas del video}\begin{itemize}}{\end{itemize}}
\newcommand{\errata}[3]{%
  \edef\errurl{https://www.youtube.com/watch?v=\videoid\string&t=#1}%
  \item[\textbf{\expandafter\href\expandafter{\errurl}{#2}}] #3}
.
% Compilar con LuaLaTeX:  lualatex nombre-de-la-nota.tex
\usepackage[spanish,provide=*]{babel}
\usepackage[T1]{fontenc}
\usepackage{lmodern}
\usepackage[a4paper,margin=2.6cm]{geometry}
\usepackage{amsmath,amssymb,amsthm}
\usepackage{xcolor}
\usepackage{tikz}
\usetikzlibrary{arrows.meta,babel}
\usepackage[colorlinks=true,urlcolor=blue!50!black,linkcolor=blue!50!black]{hyperref}

\hyphenpenalty=10000 \exhyphenpenalty=10000 \emergencystretch=3em

% Colores de las figuras (los mismos que usa la página)
\definecolor{extremo}{RGB}{47,107,79}
\definecolor{azul}{RGB}{44,74,154}

% Definiciones numeradas por unidad
\theoremstyle{definition}
\newtheorem{defn}{Definición}
% Numeración por unidad: \unidad{u}{n} (después de % Preámbulo compartido de las notas del curso de topología algebraica.
% Cada nota empieza con \documentclass[11pt]{article}, define \videoid{...}
% (el identificador del video de YouTube) y luego hace \input{preambulo}.
% Compilar con LuaLaTeX:  lualatex nombre-de-la-nota.tex
\usepackage[spanish,provide=*]{babel}
\usepackage[T1]{fontenc}
\usepackage{lmodern}
\usepackage[a4paper,margin=2.6cm]{geometry}
\usepackage{amsmath,amssymb,amsthm}
\usepackage{xcolor}
\usepackage{tikz}
\usetikzlibrary{arrows.meta,babel}
\usepackage[colorlinks=true,urlcolor=blue!50!black,linkcolor=blue!50!black]{hyperref}

\hyphenpenalty=10000 \exhyphenpenalty=10000 \emergencystretch=3em

% Colores de las figuras (los mismos que usa la página)
\definecolor{extremo}{RGB}{47,107,79}
\definecolor{azul}{RGB}{44,74,154}

% Definiciones numeradas por unidad
\theoremstyle{definition}
\newtheorem{defn}{Definición}
% Numeración por unidad: \unidad{u}{n} (después de \input{preambulo}) hace que las
% definiciones de la nota se llamen u.n, u.(n+1), ...; n es el primer número
% libre de esa unidad (ver la tabla de CONVENCIONES.md).
\newcommand{\unidad}[2]{\renewcommand{\thedefn}{#1.\arabic{defn}}\setcounter{defn}{\numexpr#2-1\relax}}

% Observaciones y erratas del video. Usar SOLO si hay algo que señalar:
%   \begin{erratas}
%     \errata{203}{3:23}{Se dice ..., pero debe decir ...}
%   \end{erratas}
% El primer argumento son los segundos (enlaza al minuto exacto del video),
% el segundo es la marca visible, el tercero es el texto.
\newenvironment{erratas}{\section*{Observaciones y erratas del video}\begin{itemize}}{\end{itemize}}
\newcommand{\errata}[3]{%
  \edef\errurl{https://www.youtube.com/watch?v=\videoid\string&t=#1}%
  \item[\textbf{\expandafter\href\expandafter{\errurl}{#2}}] #3}
) hace que las
% definiciones de la nota se llamen u.n, u.(n+1), ...; n es el primer número
% libre de esa unidad (ver la tabla de CONVENCIONES.md).
\newcommand{\unidad}[2]{\renewcommand{\thedefn}{#1.\arabic{defn}}\setcounter{defn}{\numexpr#2-1\relax}}

% Observaciones y erratas del video. Usar SOLO si hay algo que señalar:
%   \begin{erratas}
%     \errata{203}{3:23}{Se dice ..., pero debe decir ...}
%   \end{erratas}
% El primer argumento son los segundos (enlaza al minuto exacto del video),
% el segundo es la marca visible, el tercero es el texto.
\newenvironment{erratas}{\section*{Observaciones y erratas del video}\begin{itemize}}{\end{itemize}}
\newcommand{\errata}[3]{%
  \edef\errurl{https://www.youtube.com/watch?v=\videoid\string&t=#1}%
  \item[\textbf{\expandafter\href\expandafter{\errurl}{#2}}] #3}
) hace que las
% definiciones de la nota se llamen u.n, u.(n+1), ...; n es el primer número
% libre de esa unidad (ver la tabla de CONVENCIONES.md).
\newcommand{\unidad}[2]{\renewcommand{\thedefn}{#1.\arabic{defn}}\setcounter{defn}{\numexpr#2-1\relax}}

% Observaciones y erratas del video. Usar SOLO si hay algo que señalar:
%   \begin{erratas}
%     \errata{203}{3:23}{Se dice ..., pero debe decir ...}
%   \end{erratas}
% El primer argumento son los segundos (enlaza al minuto exacto del video),
% el segundo es la marca visible, el tercero es el texto.
\newenvironment{erratas}{\section*{Observaciones y erratas del video}\begin{itemize}}{\end{itemize}}
\newcommand{\errata}[3]{%
  \edef\errurl{https://www.youtube.com/watch?v=\videoid\string&t=#1}%
  \item[\textbf{\expandafter\href\expandafter{\errurl}{#2}}] #3}

\unidad{1}{23} % unidad 1; primer enunciado: 1.23

\title{Aplicaciones: el teorema de Borsuk--Ulam en dimensión 2}
\author{Luis Jorge Sánchez Saldaña\\[2pt]\small Facultad de Ciencias, UNAM}
\date{}

\begin{document}
\maketitle

\noindent\small Notas del video
\href{https://www.youtube.com/watch?v=HlZCNLpj8qw}{\guillemotleft{}El teorema de Borsuk--Ulam en dimensión 2\guillemotright{}}
(curso de Introducción a la Topología Algebraica). Tercera aplicación del cálculo $\pi_1(S^1)\cong\mathbb{Z}$: para toda función continua de la esfera $S^2$ en el plano hay dos puntos antípodas con la misma imagen.
\normalsize

\section*{Enunciado}

\begin{teo}[Teorema de Borsuk--Ulam en dimensión 2]
Para toda función continua $f\colon S^2\to\mathbb{R}^2$ existe $x\in S^2$ tal que $f(x)=f(-x)$. Es decir, existe un par de puntos antípodas de $S^2$ con exactamente la misma imagen bajo $f$.
\end{teo}

El enunciado no es nada obvio. Un ejemplo muy rápido de cómo aplicarlo: se piensa que la Tierra es una esfera y que a cada punto se le asocian su temperatura y su presión atmosférica; uno puede convencerse de que esto es una función continua con valores en $\mathbb{R}^2$. El teorema dice entonces que siempre hay dos puntos antípodas del planeta que tienen la misma temperatura y la misma presión atmosférica, lo cual es sorprendente (bajo el supuesto, claro, de que la Tierra no es plana).

\section*{La demostración}

\begin{proof}
La demostración es por contradicción. Supongamos que existe $f\colon S^2\to\mathbb{R}^2$ continua tal que $f(x)\neq f(-x)$ para todo $x\in S^2$, es decir, dos puntos antípodas cualesquiera tienen imágenes distintas.

\medskip\noindent\textit{La función $g\colon S^2\to S^1$.}
Bajo esta suposición se define
\[
  g(x)=\frac{f(x)-f(-x)}{\lVert f(x)-f(-x)\rVert}.
\]
Está bien definida porque el denominador nunca se anula, y se divide entre la norma del numerador para que $g$ caiga en $S^1$. Además $g$ es continua.

\medskip\noindent\textit{La idea.}
Se define un lazo $h$ en $S^1$ que, al mismo tiempo, será contraíble y no contraíble; esa es la contradicción. Falta entonces definir el lazo $h$ y demostrar las dos cosas.

\medskip\noindent\textit{El lazo $h$.}
Primero se define un lazo $\eta$ en $S^2$ que da una vuelta al ecuador (figura):
\[
  \eta(s)=(\cos 2\pi s,\ \sin 2\pi s,\ 0),\qquad s\in I.
\]
(La fórmula no es lo importante; sirve para tener claro cómo está definido $\eta$.) El lazo $h$ es el que se obtiene empujando $\eta$ a $S^1$ con $g$:
\[
  h=g\circ\eta\colon I\to S^1.
\]

\begin{center}
\begin{tikzpicture}[x=0.021cm,y=-0.021cm,line cap=round]
  \draw[thick] (240,150) circle (110);
  \draw[thick,azul,dashed] (130.0,150.0) -- (130.2,148.4) -- (130.6,146.9) -- (131.4,145.3) -- (132.4,143.8) -- (133.7,142.2) -- (135.4,140.7) -- (137.3,139.2) -- (139.5,137.8) -- (142.0,136.4) -- (144.7,135.0) -- (147.7,133.7) -- (151.0,132.4) -- (154.5,131.1) -- (158.3,129.9) -- (162.2,128.8) -- (166.4,127.7) -- (170.8,126.7) -- (175.3,125.7) -- (180.1,124.8) -- (185.0,124.0) -- (190.1,123.3) -- (195.3,122.6) -- (200.6,122.0) -- (206.0,121.5) -- (211.5,121.0) -- (217.1,120.7) -- (222.8,120.4) -- (228.5,120.2) -- (234.2,120.0) -- (240.0,120.0) -- (245.8,120.0) -- (251.5,120.2) -- (257.2,120.4) -- (262.9,120.7) -- (268.5,121.0) -- (274.0,121.5) -- (279.4,122.0) -- (284.7,122.6) -- (289.9,123.3) -- (295.0,124.0) -- (299.9,124.8) -- (304.7,125.7) -- (309.2,126.7) -- (313.6,127.7) -- (317.8,128.8) -- (321.7,129.9) -- (325.5,131.1) -- (329.0,132.4) -- (332.3,133.7) -- (335.3,135.0) -- (338.0,136.4) -- (340.5,137.8) -- (342.7,139.2) -- (344.6,140.7) -- (346.3,142.2) -- (347.6,143.8) -- (348.6,145.3) -- (349.4,146.9) -- (349.8,148.4) -- (350.0,150.0);
  \draw[very thick,azul] (350.0,150.0) -- (349.8,151.6) -- (349.4,153.1) -- (348.6,154.7) -- (347.6,156.2) -- (346.3,157.8) -- (344.6,159.3) -- (342.7,160.8) -- (340.5,162.2) -- (338.0,163.6) -- (335.3,165.0) -- (332.3,166.3) -- (329.0,167.6) -- (325.5,168.9) -- (321.7,170.1) -- (317.8,171.2) -- (313.6,172.3) -- (309.2,173.3) -- (304.7,174.3) -- (299.9,175.2) -- (295.0,176.0) -- (289.9,176.7) -- (284.7,177.4) -- (279.4,178.0) -- (274.0,178.5) -- (268.5,179.0) -- (262.9,179.3) -- (257.2,179.6) -- (251.5,179.8) -- (245.8,180.0) -- (240.0,180.0) -- (234.2,180.0) -- (228.5,179.8) -- (222.8,179.6) -- (217.1,179.3) -- (211.5,179.0) -- (206.0,178.5) -- (200.6,178.0) -- (195.3,177.4) -- (190.1,176.7) -- (185.0,176.0) -- (180.1,175.2) -- (175.3,174.3) -- (170.8,173.3) -- (166.4,172.3) -- (162.2,171.2) -- (158.3,170.1) -- (154.5,168.9) -- (151.0,167.6) -- (147.7,166.3) -- (144.7,165.0) -- (142.0,163.6) -- (139.5,162.2) -- (137.3,160.8) -- (135.4,159.3) -- (133.7,157.8) -- (132.4,156.2) -- (131.4,154.7) -- (130.6,153.1) -- (130.2,151.6) -- (130.0,150.0);
  \draw[thick,extremo] (310.7,173.0) -- (169.3,127.0);
  \fill[extremo] (310.7,173.0) circle (4); \fill[extremo] (169.3,127.0) circle (4);
  \fill (240.0,40.0) circle (3.5);
  \node[extremo] at (318.7,197.0) {\small $\eta(s)$};
  \node[extremo] at (183.3,103.0) {\small $\eta(s+\tfrac12)$};
  \node at (240.0,22.0) {\small polo norte};
  \node[azul] at (390,160) {\small $\eta$};
  \node at (90,90) {$S^2$};
\end{tikzpicture}
\end{center}

\noindent Figura: el lazo $\eta$ (el ecuador de $S^2$) y dos puntos antípodas $\eta(s)$ y $\eta(s+\tfrac12)$.

\medskip\noindent\textbf{Primera parte: $h$ no es contraíble.}
Primero se observa que
\[
  g(-x)=-g(x)\qquad\text{para todo }x\in S^2,
\]
lo cual se ve de inmediato de la fórmula de $g$: al cambiar $x$ por $-x$ se cambia el signo del numerador y el denominador no cambia. Esta observación dice que $g$ manda puntos antípodas de $S^2$ en puntos antípodas de $S^1$.

Como $\eta(s)$ y $\eta(s+\tfrac12)$ son puntos antípodas (figura) y $h=g\circ\eta$, se obtiene
\begin{equation}\label{eq:antipoda}
  h\bigl(s+\tfrac12\bigr)=-h(s)\qquad\text{para todo }s\in\bigl[0,\tfrac12\bigr].
\end{equation}
Esta ecuación solo tiene sentido para esos valores de $s$, pues son los únicos para los que $s+\tfrac12$ sigue en el intervalo $I$.

Como al calcular $\pi_1(S^1)$, el lazo $h$ en $S^1$ se levanta a un camino $\tilde h\colon I\to\mathbb{R}$ con $\tilde h(0)=0$ y $\varphi\circ\tilde h=h$, donde $\varphi(x)=e^{2\pi i x}$ es la proyección cubriente (la espiral). La ecuación~\eqref{eq:antipoda} implica que
\begin{equation}\label{eq:q}
  \tilde h\bigl(s+\tfrac12\bigr)=\tilde h(s)+\frac{q}{2}\qquad\text{con }q\text{ un entero impar,}
\end{equation}
donde, en principio, $q$ depende de $s$. En efecto, si se toma una preimagen de $h(s)$ y una preimagen de $h(s+\tfrac12)=-h(s)$ bajo $\varphi$, estas difieren por $\tfrac12$ módulo un entero, y es de esa diferencia entera de donde sale $q$: para pasar de una preimagen a la otra hay que avanzar una cantidad entera y en algún momento dar un brinquito de $\tfrac12$. (Convencerse de esto con la imagen de la espiral es un buen ejercicio.)

\noindent\textbf{Comentario.} Como en el cálculo de $\pi_1(S^1)$, se toma el levantamiento con punto inicial $0$, lo cual supone $h(0)=1$. Si no fuera así, se puede componer $g$ con una rotación $\rho$ de $S^1$ que mande $h(0)$ a $1$; como $\rho(-z)=-\rho(z)$, la propiedad $g(-x)=-g(x)$ se conserva y nada de lo anterior cambia.

\medskip
\textit{$q$ no depende de $s$.} De la ecuación~\eqref{eq:q} se puede despejar $q=2\bigl(\tilde h(s+\tfrac12)-\tilde h(s)\bigr)$, que es una función continua de $s\in[0,\tfrac12]$ (pues $\tilde h$ lo es). Como $q$ toma valores enteros y $[0,\tfrac12]$ es conexo, $q$ es constante por continuidad. (También es un buen ejercicio.)

\medskip
\textit{Cálculo de $\tilde h(1)$.} Se evalúa la ecuación~\eqref{eq:q} en $s=0$ y en $s=\tfrac12$:
\[
  \tilde h\bigl(\tfrac12\bigr)=\tilde h(0)+\frac q2,\qquad
  \tilde h(1)=\tilde h\bigl(\tfrac12\bigr)+\frac q2 .
\]
Sustituyendo la primera en la segunda se obtiene $\tilde h(1)=\tilde h(0)+q=q$. Es decir, el levantamiento de $h$ empieza en $0$ y termina en $q$, de modo que $h$ es homotópico (relativamente a los extremos) al lazo en $S^1$ que da $q$ vueltas, y por el Teorema~1.19 su clase en $\pi_1(S^1)\cong\mathbb{Z}$ es $[h]=q$. Como $q$ es impar, $q\neq 0$; por lo tanto $h$ \emph{no es homotópico al lazo constante}. Esto demuestra la primera afirmación.

\medskip\noindent\textbf{Segunda parte: $h$ sí es contraíble.}
Ahora se demuestra que $h$ también se puede contraer al lazo constante. Recuérdese que $h$ se obtiene de dar una vuelta al ecuador con $\eta$ y empujar con $g$. El lazo $\eta$ se puede ir contrayendo al polo norte de $S^2$: se empuja hacia el polo norte (por los paralelos) y eventualmente se llega al lazo constante en el polo norte. Si cada uno de los lazos de esta deformación se empuja a $S^1$ con $g$, se obtiene una deformación de $h$ al lazo constante.

\medskip\noindent\textbf{Comentario.} Como en la demostración del teorema de Brouwer, esta homotopía deforma $\eta$ al lazo constante en el polo norte, y al hacerlo \emph{mueve el punto base}: los extremos de $\eta_t$ recorren un meridiano desde $\eta(0)=(1,0,0)$ hasta el polo norte. Entonces $g\circ\eta_t$ es una homotopía libre de $h$ a un lazo constante, no relativa a los extremos. Esto basta: si $c(t)=g(\eta_t(0))$ es el camino que recorre el punto base, una homotopía libre de lazos da $[h]=\beta_c([k])$, donde $k$ es el lazo constante final y $\beta_c$ es el isomorfismo de cambio de punto base (Definición~1.13, Teorema~1.12), y como $[k]$ es el neutro, también lo es $[h]$. Así se completa la segunda afirmación.

\medskip\noindent\textit{La contradicción.}
Hemos demostrado que $h$ es no contraíble (primera parte) y contraíble (segunda parte), lo cual es una contradicción. Por lo tanto, no puede existir una $f$ con $f(x)\neq f(-x)$ para todo $x$: existe $x\in S^2$ con $f(x)=f(-x)$, y esto demuestra el teorema.
\end{proof}

\begin{erratas}
  \errata{643}{10:43}{Se contrae $\eta$ al polo norte y se concluye que $h$ es homotópico al lazo constante; pero esa homotopía mueve los extremos del lazo, así que no es relativa a los extremos y no basta para contradecir $[h]=q\neq 0$ en $\pi_1(S^1)$. Basta observar que una homotopía libre de $h$ a un lazo constante da $[h]=\beta_c([k])$, que es el neutro (cambio de punto base).}
\end{erratas}

\end{document}
